\documentclass[11pt, a4paper]{report}
%----------------------------------------------------------------------------------------
%	PACKAGES (宏包管理)
%----------------------------------------------------------------------------------------
% 1. 编码与语言（中文支持）
\usepackage[UTF8, scheme=plain]{ctex}
% scheme=plain 保持定理/章节等名称为英文；若需全中文改为 scheme=chinese
% 2. 页面布局与排版
\usepackage{geometry}
\geometry{
    textheight=9in,
    textwidth=6in, % 稍微加宽以便放入更好的box
    top=1in,
    bottom=1in,
    headheight=12pt,
    headsep=25pt,
    footskip=30pt
}
\usepackage{setspace}
\setstretch{1.2}
\usepackage{float}
\usepackage{enumitem} % 更好的列表控制
\usepackage{tocloft}  % 加载目录控制宏包
\renewcommand{\cftsecleader}{\cftdotfill{\cftdotsep}} % 强制让 section 后面显示点号
% 3. 数学支持
\usepackage{amsmath, amsthm, amssymb, amsfonts}
\usepackage{mathtools}
\usepackage{thmtools} % 定理环境增强
\usepackage{extarrows}
% 4. 图形与颜色
\usepackage{graphicx}
\usepackage[dvipsnames, svgnames]{xcolor} % 加载更多颜色名
\usepackage{tikz}
\usepackage{tikz-cd}
\usepackage{ytableau} % 杨表
\usetikzlibrary{decorations.pathreplacing, calc, shadows}
% 5. 漂亮的文本框
\usepackage{tcolorbox}
\tcbuselibrary{breakable, skins, theorems}
% 6. 参考文献
%\usepackage[backend=biber,style=gb7714-2015,sorting=gb7714-2015]{biblatex}
%\addbibresource{references.bib}
% 7. 交互 (最后加载)
\usepackage[hidelinks]{hyperref}
\usepackage{array}
%----------------------------------------------------------------------------------------
%	CUSTOM COLORS (自定义配色方案)
%----------------------------------------------------------------------------------------
% 定理色系 - 蓝色
\definecolor{ThmColor}{RGB}{0, 71, 171}        % 深钴蓝
\definecolor{ThmBg}{RGB}{240, 244, 255}        % 极浅蓝
% 定义色系 - 绿色
\definecolor{DefColor}{RGB}{0, 100, 0}         % 深绿
\definecolor{DefBg}{RGB}{240, 255, 240}        % 蜜瓜绿
% 引理/命题色系 - 橙色
\definecolor{LemColor}{RGB}{204, 85, 0}        % 焦橙色
\definecolor{LemBg}{RGB}{255, 248, 240}        % 极浅橙
% 示例/注记色系 - 灰色/紫色
\definecolor{ExColor}{RGB}{105, 105, 105}      % 暗灰
\definecolor{ExBg}{RGB}{250, 250, 250}         % 极浅灰
% 专题/话题色系 - 靛紫色
\definecolor{TopicColor}{RGB}{75, 0, 130}      % 靛蓝色 (Indigo)
\definecolor{TopicBg}{RGB}{248, 245, 255}      % 极浅紫
%----------------------------------------------------------------------------------------
%	THEOREM ENVIRONMENTS (环境美化核心)
%----------------------------------------------------------------------------------------
% 1. Theorem (蓝色风格)
\declaretheorem[name=Theorem, numberwithin=section]{theorem}
\tcolorboxenvironment{theorem}{
    enhanced jigsaw,
    colframe=ThmColor,       % 边框颜色
    colback=ThmBg,           % 背景颜色
    coltitle=white,          
    boxrule=0.5pt,           % 边框粗细
    leftrule=2pt,            % 左边框加粗
    sharp corners=downhill,  % 样式
    breakable,               % 允许跨页
    drop shadow=black!5!white % 轻微阴影
}
% Corollary 使用与 Theorem 相同的风格，但通常紧随其后
\declaretheorem[name=Corollary, numberlike=theorem]{corollary}
\tcolorboxenvironment{corollary}{
    enhanced jigsaw,
    colframe=ThmColor!80!white, 
    colback=ThmBg,
    boxrule=0pt, leftrule=2pt, arc=0pt, outer arc=0pt,
    breakable
}
% 2. Definition (绿色风格)
\declaretheorem[name=Definition, numberlike=theorem]{definition}
\tcolorboxenvironment{definition}{
    enhanced jigsaw,
    colframe=DefColor,
    colback=DefBg,
    boxrule=0.5pt, leftrule=2pt,
    sharp corners=downhill,
    breakable,
    drop shadow=black!5!white
}
% 3. Lemma & Proposition (橙色风格)
\declaretheorem[name=Lemma, numberlike=theorem]{lemma}
\declaretheorem[name=Proposition, numberlike=theorem]{proposition}
\tcolorboxenvironment{lemma}{
    enhanced jigsaw,
    colframe=LemColor,
    colback=LemBg,
    boxrule=0.5pt, leftrule=2pt,
    breakable
}
\tcolorboxenvironment{proposition}{
    enhanced jigsaw,
    colframe=LemColor,
    colback=LemBg,
    boxrule=0.5pt, leftrule=2pt,
    breakable
}
% 4. Example & Remark (简约风格)
\declaretheorem[name=Example, numberlike=theorem]{example}
\declaretheorem[name=Remark, numberlike=theorem]{remark}
\declaretheorem[name=Exercise, numberlike=theorem]{exercise}
\tcolorboxenvironment{example}{
    enhanced jigsaw,
    colframe=ExColor,
    colback=ExBg,
    leftrule=2pt, boxrule=0pt,
    frame hidden, % 隐藏其他三边
    breakable
}
% 5. Topic (靛紫色风格，带标题支持)
\declaretheorem[name=Topic, numberlike=theorem]{topic}
\tcolorboxenvironment{topic}{
    enhanced jigsaw,
    breakable,
    colframe=TopicColor,        % 边框为靛紫色
    colback=TopicBg,           % 背景为极浅紫
    boxrule=0.5pt,             % 细边框
    leftrule=2pt,              % 左侧加粗线
    sharp corners=downhill,    % 与你其他环境一致的切角样式
    drop shadow=black!5!white, % 轻微阴影
    before skip=15pt,
    after skip=15pt
}
%----------------------------------------------------------------------------------------
%	MATH MACROS (保留你的常用定义)
%----------------------------------------------------------------------------------------
\newcommand{\g}{\mathfrak{g}}  
\newcommand{\der}{\mathrm{d}} 
\newcommand{\h}{\mathfrak{h}} 
\newcommand{\n}{\mathfrak{n}}  
\newcommand{\U}{\mathcal{U}}  
\newcommand{\C}{\mathbb{C}} 
\newcommand{\R}{\mathbb{R}} 
\newcommand{\borel}{\mathfrak{b}}
\newcommand{\Ind}{\mathrm{Ind}} 
\newcommand{\Hom}{\mathrm{Hom}}  
\newcommand{\Ext}{\mathrm{Ext}}
\newcommand{\HRule}[1]{\rule{\linewidth}{#1}}
%----------------------------------------------------------------------------------------
%	TITLE PAGE
%----------------------------------------------------------------------------------------
\begin{document}
\begin{titlepage}
    \centering
    \vspace*{2cm}
    
    \textsc{\Large Course Notes}\\[0.5cm]

    \HRule{1.5pt} \\[0.5cm]
    { \huge \bfseries Lie Theory }\\[0.4cm] 
    \HRule{1.5pt} \\[1.5cm]
    
    \vspace{3cm}

    % 人员信息区域：垂直排列并区分字号
    \begin{spacing}{1.8}
        {\Large \textbf{Instructor:}} \\
        {\Large Prof. Efilm \textsc{Zelmanov}} \\[0.8cm]
        
        {\normalsize \textit{Student:}} \\
        {\normalsize Leo}
    \end{spacing}

    \vfill
    {\large \today}
\end{titlepage}
\newpage
\tableofcontents
\newpage


%----------------------------------------------------------------------------------------
%    CONTENT STRUCTURE TEMPLATE
%----------------------------------------------------------------------------------------
\chapter{Lecture 1}

\section{Introduction}

Lie algebras were originally introduced by Sophus Lie in the 1870s to study the concept of continuous transformation groups. Lie's motivation was to develop a theory for differential equations analogous to Galois theory for algebraic equations. He discovered that the infinitesimal transformations of a continuous group form a linear space equipped with a bracket operation, now known as a Lie algebra. This linearization allows complex global geometric problems to be translated into algebraic problems, making them more tractable.

Let $ F $ be a field (mostly $ \operatorname{char}F =0$, even $ F=\mathbb{C} $). 

\begin{definition} 
A linear algebra $ A $ is a vector space over $ F $, with a binary bilinear operation:
\begin{itemize}
    \item binary operation: $ A \times A \to A $, $ (a,b) \mapsto ab $.
    \item bilinear: this operation is linear in each variable. In orther word, $ \forall \alpha \in F $ and $ a,a',a'',b,b',b'' \in A $, we have:
    \begin{align*}
    (\alpha a)b=a(\alpha b) = \alpha (ab) \\
    (a'+a'')b = a'b+a''b \\
    a(b'+b'')=ab'+ab''
    \end{align*}
\end{itemize}
The algebra $ A $ is associative if $ \forall a,b,c \in A $, $ (ab)c=a(bc) $.
\end{definition}

\begin{example} 
$ M_n(F) $, the algebra of $ n \times n $ matrices over $ F $.
\end{example}

\begin{example} 
Let $ V $ be a vector space, $ \operatorname{End}_F (V) =\{ \text{linear transformations } \varphi: V\to V \}$ is an associative algebra with the composition operation$ (\phi \psi) (v)=\phi(\psi(v))$. 

If $ \dim V < \infty $, when we fix a basis, then we fet a matrix representation of $ \varphi $ in our fix basis.
\end{example}

The most imprtant thing everyone shoul know is Wedderburn's theorem.

\begin{theorem}[Wedderburn-Artin]
    Let $A$ be a semisimple Artinian ring (with 1). Then $A$ is a direct sum of matrix rings over division rings. 
    
    Specifically, if ${}_A A \simeq S_1^{n_1} \oplus \dots \oplus S_r^{n_r}$ where $S_1, \dots, S_r$ are non-isomorphic simple modules occurring with multiplicities $n_1, \dots, n_r$, then
    \[
    A \simeq M_{n_1}(D_1) \oplus \dots \oplus M_{n_r}(D_r) \quad 
    \]
    where $D_i = \operatorname{End}_A(S_i)^{op}$. This is a ring isomorphism, and each $M_{n_i}(D_i)$ is a 2-sided ideal.
    \end{theorem}
    
    \begin{proof}
    First, since $A$ is a semisimple ring, its left regular module ${}_A A$ can be decomposed into a direct sum of simple left $A$-modules:
    \[
    {}_A A \simeq S_1^{n_1} \oplus \dots \oplus S_r^{n_r} \quad 
    \]
    where $S_1, \dots, S_r$ are pairwise non-isomorphic simple modules.
    
    By Schur's lemma, the endomorphism ring of a simple module is a division ring, so $D_i = \operatorname{End}_A(S_i)^{op}$ is a division ring.
    
    Next, we compute the endomorphism ring of the left regular module $\operatorname{End}_A({}_A A)$. On one hand, for any ring $A$ with 1, we have the ring isomorphism $\operatorname{End}_A({}_A A) \simeq A^{op}$. On the other hand, using our direct sum decomposition:
    \[
    \operatorname{End}_A({}_A A) \simeq \operatorname{End}_A(S_1^{n_1}) \oplus \dots \oplus \operatorname{End}_A(S_r^{n_r}) \quad 
    \]
    This decomposition holds because $\operatorname{Hom}_A(S_i^{n_i}, S_j^{n_j}) = 0$ whenever $i \neq j$. 
    
    For each component in the direct sum, we can rewrite the endomorphism ring using matrices over $D_i$:
    \[
    \operatorname{End}_A(S_i^{n_i}) \simeq M_{n_i}(D_i)^{op} \simeq M_{n_i}(D_i^{op}) \quad 
    \]
    Substituting this back into our equation for $A^{op}$, we obtain:
    \[
    A^{op} \simeq M_{n_1}(D_1^{op}) \oplus \dots \oplus M_{n_r}(D_r^{op})
    \]
    
    To recover the ring $A$, we take the opposite ring on both sides. Recall the property that $M_n(R)^{op} \simeq M_n(R^{op})$. Applying this to each block yields:
    \[
    A \simeq (A^{op})^{op} \simeq \left( \bigoplus_{i=1}^r M_{n_i}(D_i^{op}) \right)^{op} \simeq \bigoplus_{i=1}^r M_{n_i}(D_i)
    \]
    Thus, the isomorphism is established.
\end{proof}
    
    \begin{remark}
    If the base field $F$ is algebraically closed, every finite-dimensional
    division algebra over $F$ is equal to $F$. Hence the theorem simplifies to
    \[
    A \cong \prod_{i=1}^{r} M_{n_i}(F).
    \]
    \end{remark}
    
    \begin{example}
    Let $G$ be a finite group and let $\mathbb{C}[G]$ be its group algebra.
    By Maschke's theorem $\mathbb{C}[G]$ is semisimple. Therefore
    \[
    \mathbb{C}[G] \cong \prod_{i=1}^{r} M_{n_i}(\mathbb{C}),
    \]
    where $n_i$ are the dimensions of the irreducible complex representations
    of $G$.
    \end{example}

\begin{definition} 
Suppose $ F $ is algebraically closed. An algebra $ A $ is simple if 
\begin{enumerate}
    \item $ \nexists \  I \triangleleft A $ and $ I \neq 0 $. 
    \item $ A\cdot  A \neq (0)$.   
\end{enumerate}
\end{definition}

\begin{corollary}
Any simple finite dimensional associative algebra is isomorphic to $ M_n(F) $ for some $ n $.
\end{corollary}

\begin{definition}[Jacobson Radical]
    Let $A$ be an associative algebra with identity over a field $F$.
    The \emph{Jacobson radical} of $A$, denoted by $\operatorname{Rad}(A)$
    or $J(A)$, is defined as the intersection of all maximal left ideals of $A$:
    \[
    \operatorname{Rad}(A)
    =
    \bigcap_{M \text{ maximal left ideal of } A} M.
    \]
    \end{definition}
    
    
    \begin{proposition}
    Let $A$ be an associative algebra with identity. Then
    \[
    \operatorname{Rad}(A)
    =
    \bigcap_{S \text{ simple left } A\text{-module}}
    \operatorname{Ann}_A(S),
    \]
    where
    \[
    \operatorname{Ann}_A(S)
    =
    \{a\in A \mid aS=0\}
    \]
    is the annihilator of $S$ in $A$.
    \end{proposition}
    
    
    \begin{proposition}[Equivalent Characterization]
    Let $A$ be an associative algebra with identity.
    Then an element $a\in A$ lies in $\operatorname{Rad}(A)$
    if and only if for every $x\in A$ the element
    \[
    1-ax
    \]
    is invertible in $A$.
    \end{proposition}
    
    
    \begin{proposition}
    The Jacobson radical $\operatorname{Rad}(A)$ is a two-sided ideal of $A$.
    \end{proposition}
    
    
    \begin{definition}[Semisimple Algebra]
    An associative algebra $A$ is called \emph{semisimple} if
    \[
    \operatorname{Rad}(A)=0.
    \]
    \end{definition}

\section{Where do Lie algebras come from?}

\begin{definition} 
A Lie algebra is a vector space $ L $ over a field $ F $, with a binary operation $ [a,b] \in L $, such that
\begin{enumerate}
    \item $ [a,b] $ is bilinear. 
    \item $ [a,b] = -[b,a] $.
    \item $ [a,[b,c]] + [c,[a,b]] + [b,[c,a]] = 0 $, this is called Jacobi identity. 
\end{enumerate}
\end{definition}

\subsection{They come from associative (central simple) algebras}

Let $ A $ be an associative algebra. $ \forall a,b \in A $, define $ [a,b]=ab-ba $. Check that $ (A,[ \ ,\  ]) $ is a Lie algebra:

\begin{proof}
    Bilinearity and skew-symmetry are clear. For Jacobi identity:
    \begin{align*}
        [a,[b,c]] & + [b,[c,a]] + [c,[a,b]] \\ 
        &= a(bc-cb) - (bc-cb)a + b(ca-ac) - (ca-ac)b + c(ab-ba) - (ab-ba)c \\
        &= abc - acb - bca + cba + bca - bac - cab + acb + cab - cba - abc + bac \\
        &= 0.
    \end{align*}
\end{proof}

We denote this new algebra $ (A,[ \ ,\  ]) $ by $ A^{(-)} $. 

Suppose that $ L \subseteq A $ and $ [L,L] \subseteq L $. Then $ L \le A^{(-)} $。

\begin{example} 
$ M_n(F)^{(-)}=\mathfrak{gl}(n) $. We call it the general linear Lie algebra.

$ [L,L]=\{ \sum_{i} [a_i,b_i] \} \trianglelefteq L $. Clearly $ [\mathfrak{gl}(n),\mathfrak{gl}(n)] \neq \mathfrak{gl}(n) $ and for $ n \ge 2 $, $ [\mathfrak{gl}(n),\mathfrak{gl}(n)]\neq 0 $.    
\end{example}

\begin{example} 
Focus on $ [\mathfrak{gl}(n),\mathfrak{gl}(n)] $, all matrices are of trace 0. This is called the special linear Lie algebra, denoted by $ \mathfrak{sl}(n) $.
\end{example}

\begin{definition} 
$ Z $ is the center of Lie algrbra $ L $ if $ Z=\{ z \in L | [z,L]=0 \} $.   
\end{definition}

\subsection{They come from involutions of central simple algebras}

We can also find Lie algebras from the involution of a central simple algebra. Let $ A $ be central simple algebra and $ *: A \to A $ be a linear transformation. $ * $ is called a an involution if $ (a^*)^* =a $ and $ (ab)^*=b^*a^* $, $ \forall a,b \in A $.

\begin{proposition}
    Let $A$ be a central simple algebra over a field $F$ equipped with an involution
    $\sigma : A \to A$. Then the set
    \[
    L=\{x\in A \mid \sigma(x)=-x\}
    \]
    forms a Lie algebra under the commutator bracket
    \[
    [x,y]=xy-yx .
    \]
    \end{proposition}
    
    \begin{proof}
    First observe that any associative algebra $A$ becomes a Lie algebra
    with the commutator bracket
    \[
    [x,y]=xy-yx .
    \]
    
    Let
    \[
    L=\{x\in A\mid \sigma(x)=-x\}.
    \]
    We show that $L$ is closed under the Lie bracket.
    
    Take $x,y\in L$. Then
    \[
    \sigma(x)=-x,\qquad \sigma(y)=-y.
    \]
    
    Since $\sigma$ is an involution, it satisfies
    \[
    \sigma(xy)=\sigma(y)\sigma(x).
    \]
    
    Hence
    \[
    \sigma([x,y])
    =\sigma(xy-yx)
    =\sigma(y)\sigma(x)-\sigma(x)\sigma(y).
    \]
    
    Substituting $\sigma(x)=-x$ and $\sigma(y)=-y$, we obtain
    \[
    \sigma([x,y])
    =(-y)(-x)-(-x)(-y)
    =yx-xy
    =-[x,y].
    \]
    
    Thus $[x,y]\in L$. Therefore $L$ is closed under the commutator bracket.
    
    Since the commutator bracket in any associative algebra satisfies
    bilinearity, antisymmetry, and the Jacobi identity,
    it follows that $L$ is a Lie algebra.
    \end{proof}







\begin{example}
 The Quatanion algebra is an associative algebra with the involution $ a \mapsto a^* $. It is 4 dimensional central simple algebra with standard basis $ 1,i,j,k $, where the relations are $ i^2 = j^2 = k^2  = -1, ij=k, jk=i, ki=j $ as the following graph:
 \[
    \begin{tikzcd}
        i \arrow[r] & j \arrow[d]  \\
                    & k \arrow[lu]
        \end{tikzcd}
 \]
  We use $ \mathbb{H} $ to denote the Quatanion algebra when we choose $ \mathbb{R} $ as the base field. i.e. $ \mathbb{H}=\{ a_0 1 + a_1 i + a_2 j + a_3 k | a_0,a_1,a_2,a_3 \in \mathbb{R} \} $. $ \forall x \in \mathbb{H} $ where $ x=a_0 1 + a_1 i + a_2 j + a_3 k $, $ * $ is defined as $ x^* = a_0 1 - a_1 i - a_2 j - a_3 k $ and $ x^{-1}= \frac{1}{a_0^2 + a_1^2 + a_2^2 + a_3^2} (a_0 1 - a_1 i - a_2 j - a_3 k) $. Hence $ \mathbb{H} $ is a division algebra. 
 
 If the base field is algebraically closed, $F =\overline{F}$. We can denote the Quatanion algebra by $ 2 \times 2 $ matrices 
 \[
 \begin{pmatrix}
    \alpha & \beta \\
    \gamma & \xi \\
 \end{pmatrix}
 \]    
Where $ \alpha, \beta, \gamma, \xi \in F $. In this case it is no longer a division algebra ($a_0^2 + a_1^2 + a_2^2 + a_3^2$ may equal to 0), but just $ M_2(F) $. And in this case the involution $ * $ is defined as:
\[
\begin{pmatrix}
    \alpha & \beta \\
    \gamma & \xi \\
 \end{pmatrix}
 \mapsto
 \begin{pmatrix}
    \xi & -\beta \\
    -\gamma & \alpha \\
 \end{pmatrix}
\]   
Which is the Simplectic involution.
\end{example}

\begin{example} 

\end{example}















\chapter{Lecture 2}

Recall last time: 

1. Matrix Lie algebras, e.g. $ \mathfrak{sl}(n) $. Consider involution $ * : M_n(F) \to M_n(F) $ skew-symmetric matrices.

2. Derivation, $ \operatorname{Der}(A) $. e.g. $ A=F[t_1,\ldots ,t_n] $, $ \mathrm{d}=\sum_{min}^{max} f_i() \partial $

A non-associative algebra, $ \mathrm{d}: x \mapsto [a,x] $ is a derivation. 

Let $ L $ a Lie algebra, inner derivation

Q: $ \mathrm{d} $ a derivation and $ \operatorname{exp}(\mathrm{d}) $. It need convergence.

\begin{example} 
$ \mathrm{d}=t_1t_2 \frac{\partial}{\partial t_1} $.  
\end{example}

3. Brackets: 

\section{Engel Theorem}

$ L^n $: span of products of elements of length $ n $. Recall $ \mathrm{ad}([a,b])=[\mathrm{ad}(a),\mathrm{ad}(b)] $. $ L \to \operatorname{End}_F(L)^{(-)} , a \mapsto \mathrm{ad}(a)$. 

$ [I,J] $ is again an ideal.

Decending chain:
\[
L=L^1 \supseteq L^2 \supseteq \ldots
\]
L is nilpotent if....

eg: strict upper matrix

Another decending chain: $ L^{[1]}=[L,L] $, $ L^{[n+1]}=[L^{[n]},L^{[n]}] $  

L is solvable if .....

e.g. 2 dimentional sovable non nilpotent Lie alg.

$ A $ associa. $ L \subseteq A^{(-)} $. Suppose $ A = <L> $ as an associa alge. Then A is an enveloping algebra of $ L $.

\begin{theorem} 
Let $ A $ be a finite dimensional associative algebra, $ L\subseteq A^{(-)} $, $ A = <L> $. $ \exists n \ge 1 $, $ \forall a \in L $, $ a^n=0 $. Then $ A $ is nilpotent.       
\end{theorem}

We say that a subset $S \subseteq L$ is a Lie subset if $ \forall a,b \in S $, $ [a,b] \in S $.

\begin{theorem} 
Let $ A $ be a finite dimensional associative algebra, $ L\subseteq A^{(-)} $, $ A = <L> $. Let $ S $ be a Lie subset of $ L $. $ L = \mathrm{span}(S) $, $ \forall a \in S $, $ a^n=0 $.    
\end{theorem}

Choose a maximal Lie subset $ S_1 \subset S $ such that the associative $ <S_1> $ is nilpotent. So $ \exists n $ such that $ S_1 \ldots S_1 =(0) $ (n times)   (existance?)

claim $ \mathrm{ad}(S_1) \ldots \mathrm{ad}(S_n)=(0) $.

  




\end{document}
